% ---------
% --------
%!TEX TS-program = pdflatex
%!TEX encoding = UTF-8 Unicode
%!TEX spellcheck = English (Aspell)

\documentclass[11pt]{article}
\usepackage{../jeffe,../handout,graphicx}

\begin{document}

\begin{center}
\Large\textbf{CSCI 432/532, Spring 2026}%
\\
\LARGE\textbf{Final Project}%
\end{center}

\bigskip
\hrule
\bigskip

The final project for this course is to design and deliver a lecture covering an algorithms topic that we did not
cover this semester. In groups of 3, you will select a topic, use outside resources
(I can provide some suggestions) to learn about it, and prepare a 50 minute lecture
to teach the class about the topic in the context of what we have learned in this course.

\section{Project Schedule \& Deliverables}
\begin{itemize}
    \item \textbf{Friday, March 27, by 10am: group formation}. No deliverable, but make sure I know your group by this date.
    Every group must have exactly 3 members: one graduate student and two undergraduates.
    \item \textbf{Friday, April 3, by 5pm: topic selection}: Submit a short writeup on Canvas (only one group member will need to submit) giving a \emph{formal description of the problem or topic} and
    at least two quality resources (textbooks, collections of lecture videos, etc.) that you will use to prepare your lecture.
    \item \textbf{Friday, April 15, in class: teaser presentation}: Give a 10 minute presentation to the class
    introducing your topic: the basic problem definition, why you chose it, and what
    you still are working to understand about it. And get us excited to learn more about it!
    \item \textbf{April 29--May 4 in class and finals period}: Deliver a 50 minute lecture lecture to the class on your topic. You must also
    prepare a handout for the class providing any formal definitions, theorems,
    and proofs that you cover in your lecture.
\end{itemize}

\section{Grading}
Your total project grade will be based on the following. Note that 85 points are 
earned with the entire group, and 15 points are earned individually by participating
and providing feedback on other groups' presentations.
\begin{itemize}
    \item 0 points: group formation.
    \item 5 points: topic selection deliverable.
    \item 10 points: teaser presentation.
    \item 10 points: handout for final presentation.
    \item 55 points: final presentation. 
    \item 5 points: self and peer evaluation of group members.
    \item 15 points (individual): audience participation during and providing feedback on other groups' presentations.
\end{itemize}

\section{Potential Topics}

Here are some topics to consider. Feel free to choose a topic outside of this list. Note that some topics may
be too large for a single lecture, so you may need to focus on a specific aspect
of the topic. Additionally, some topiocs may require more background than we have covered in class,
but they may still be worth taking on if you are excited about them! I am happy to advise on
the choice of a topic.
\begin{itemize}
    \item Advanced computer science theory topics: approximation algorithms,
    fixed paramter tractability, fine-grained complexity,
    proof complexity, space complexity, attacks on P vs. NP,
    etc.
    \item Typical advanced algorithms topics that we don't get to: network
    flows, linear programming, randomized algorithms, distributed algorithms,
    etc.
    \item Advanced data structures: range queries, hashing and sketching, bloom
    filters, tries and suffix trees, etc.
    \item Algorithmic components of training deep neural networks: back
    propogation and gradient descent
    \item Quantum speedups for machine learning algorithms (quantum annealing,
    HHL algorithm, quantum kernel method, quantum PCA, etc.)
    \item More advanced cryptography topics: zero-knowledge proofs, blockchains, etc.
    \item Practical solutions to NP-hard problems: integer linear programming,
    SAT solvers, TSP solvers, etc.
\end{itemize}
 
%%%%%%%%%%%%%%%%%%%%
\end{document}

